\section{Normal forms}
\label{sec:prop-syntax-semantics-04}

\begin{definition}[Propositional literal, clause, and formula interpretation]
A \emph{propositional literal} is a propositional variable or the negation of
a propositional variable.  A \emph{clause} is a finite set of propositional
literals.  Thus clause equality is set equality: order and repeated
occurrences are not part of a clause.

To map this set object back to the formula language, put
\[
 \operatorname{code}(p_n)=2n,
 \qquad
 \operatorname{code}(\neg p_n)=2n+1.
\]
For a clause \(C\), let \(k=\abs{C}\); when \(k>0\), list its members as
\(L_1,\ldots,L_k\) in strictly increasing code order and define, by downward
recursion,
\[
 D_k\defeq L_k,
 \qquad
 D_i\defeq(L_i\lor D_{i+1})\quad(1\leq i<k).
\]
Then define
\[
 \operatorname{form}_{\CPL}(C)\defeq
 \begin{cases}
  \bottom,&C=\varnothing,\\
  D_1,&C\ne\varnothing.
 \end{cases}
\]
This is an exact metalinguistic function from clauses to
\(\Fm_{\CPL}\), not a new formula constructor.  A valuation \(v\)
\emph{satisfies} \(C\) iff
\(v\entails\operatorname{form}_{\CPL}(C)\).  In particular, the empty
clause is false under every valuation.
\end{definition}

\begin{definition}[Negation, conjunctive, and disjunctive normal form]
A formula is in \emph{negation normal form} iff it is built from literals and
\(\bottom,\topv\) using only \(\land\) and \(\lor\).  It is in
\emph{conjunctive normal form} iff it is \(\topv\), the empty conjunction,
or there are clauses \(C_1,\ldots,C_m\), with \(m>0\), such that it is
\(N_1\), where the formula sequence \(N_m,\ldots,N_1\) is defined by
\[
 N_m\defeq\operatorname{form}_{\CPL}(C_m),
 \qquad
 N_i\defeq
 \bigl(\operatorname{form}_{\CPL}(C_i)\land N_{i+1}\bigr)
 \quad(1\leq i<m).
\]
Thus a conjunctive-normal-form formula is a right-associated finite
conjunction of the formula interpretations of clauses.  In every finite
conjunction and finite disjunction used in this chapter, the empty conjunction
is \(\topv\) and the empty disjunction is \(\bottom\).  A formula is in
\emph{disjunctive normal form} iff it is a finite disjunction of finite
conjunctions of literals under these conventions.
\end{definition}

\begin{lemma}[Satisfaction of clauses and of clause conjunctions]
\label{lem:prop-clause-satisfaction}
Let \(v\) be a valuation.
\begin{enumerate}
  \item For every clause \(C\):
        \(v\entails\operatorname{form}_{\CPL}(C)\) iff
        \(v\entails L\) for at least one literal \(L\in C\).
  \item For all clauses \(C_1,\ldots,C_m\) with \(m>0\), and the formulas
        \(N_m,\ldots,N_1\) built from them as in the definition of
        conjunctive normal form: \(v\entails N_1\) iff
        \(v\entails\operatorname{form}_{\CPL}(C_j)\) for every \(j\) with
        \(1\leq j\leq m\).
\end{enumerate}
\end{lemma}
\begin{proof}
For 1, if \(C=\varnothing\), then
\(\operatorname{form}_{\CPL}(C)=\bottom\), which no valuation satisfies,
and \(C\) has no member, so both sides fail.  Otherwise let
\(k=\abs{C}>0\), list \(C\) in strictly increasing code order as
\(L_1,\ldots,L_k\), and let \(D_k,\ldots,D_1\) be the formulas of the
clause-interpretation recursion.  By downward induction along that
recursion, \(v\entails D_i\) iff \(v\entails L_j\) for some \(j\) with
\(i\leq j\leq k\): the base case \(D_k=L_k\) is immediate, and the case
\(D_i=(L_i\lor D_{i+1})\) follows from the recursive truth clause for
\(\lor\) together with the induction hypothesis.  Taking \(i=1\) and
\(\operatorname{form}_{\CPL}(C)=D_1\) proves 1.

For 2, by downward induction along the recursion defining
\(N_m,\ldots,N_1\), \(v\entails N_i\) iff
\(v\entails\operatorname{form}_{\CPL}(C_j)\) for every \(j\) with
\(i\leq j\leq m\): the base case \(N_m=\operatorname{form}_{\CPL}(C_m)\)
is immediate, and the case
\(N_i=(\operatorname{form}_{\CPL}(C_i)\land N_{i+1})\) follows from the
recursive truth clause for \(\land\) together with the induction
hypothesis.  Taking \(i=1\) proves 2.
\end{proof}

\begin{theorem}[Negation normal form]
\label{thm:p02-w02-p02-045-negation-normal-form}
Every propositional formula is logically equivalent to a formula in negation
normal form.
\end{theorem}
\begin{proof}
First eliminate every occurrence of \(\to\) using
\(\varphi\to\psi\equiv\neg\varphi\lor\psi\).  Proceed by structural
induction through the resulting formula.  Double negations are removed, and a
negation immediately above a binary connective is moved inward by
\[
 \neg(\varphi\land\psi)\equiv\neg\varphi\lor\neg\psi,
 \qquad
 \neg(\varphi\lor\psi)\equiv\neg\varphi\land\neg\psi.
\]
The base cases are variables, their negations, \(\bottom\), and \(\topv\).
Each displayed replacement decreases the number of connectives lying below a
negation, so the process terminates.  Every replacement is checked by the
recursive truth clauses and therefore preserves logical equivalence.
\end{proof}

\begin{theorem}[Conjunctive and disjunctive normal forms]
\label{thm:prop-normal-forms}
Every propositional formula is logically equivalent to a formula in
disjunctive normal form and to a formula in conjunctive normal form.
\end{theorem}
\begin{proof}
Let \(p_1,\ldots,p_n\) be the variables occurring in \(\varphi\), listed in
increasing index order.  In this
proof only, a \emph{valuation on} a finite variable set is a function from
that set to \(\B\); the recursive truth clauses apply verbatim to formulas
whose variables lie in the set, and \cref{lem:prop-coincidence} identifies
satisfaction under such a function with satisfaction under any total
valuation extending it.  For each
valuation \(v\) on these variables define
\[
  \delta_v\defeq\bigwedge_{i=1}^{n}\ell_i^v,
  \qquad
  \ell_i^v\defeq
  \begin{cases}
    p_i,&v(p_i)=\True,\\
    \neg p_i,&v(p_i)=\False.
  \end{cases}
\]
A valuation on the empty variable set is the empty function, and when
\(n=0\) the displayed conjunction is \(\topv\) by the preceding convention.
A valuation \(w\) satisfies \(\delta_v\) iff \(w(p_i)=v(p_i)\) for every
\(i\).  Let \(V_{\varphi}\) be the finite set of valuations on
\(p_1,\ldots,p_n\) that satisfy \(\varphi\).  Order these valuations
lexicographically by
\((v(p_1),\ldots,v(p_n))\), with \(\False<\True\), enumerate the resulting
order as \(v_1,\ldots,v_m\), and put
\[
  D_{\varphi}\defeq\bigvee_{j=1}^{m}\delta_{v_j}.
\]
If \(V_{\varphi}=\varnothing\), then \(m=0\) and this disjunction is
\(\bottom\).  By the
coincidence lemma, an arbitrary valuation satisfies \(D_{\varphi}\) iff
it satisfies \(\varphi\).  Hence \(D_{\varphi}\) is the required
disjunctive normal form.

For the conjunctive normal form, argue by structural induction on
\(\varphi\).  Say that a valuation \(w\) \emph{satisfies} a finite set
\(\mathcal A\) of clauses iff it satisfies every member of \(\mathcal A\).
Every valuation satisfies the empty set of clauses, and no valuation
satisfies a set that contains the empty clause.  Two facts about clause
sets are used.
\begin{enumerate}
  \item[(a)] For clauses \(A\) and \(B\), a valuation \(w\) satisfies
        \(A\cup B\) iff it satisfies \(A\) or satisfies \(B\).  By
        \cref{lem:prop-clause-satisfaction}(1), \(w\) satisfies \(A\cup B\)
        iff \(w\entails L\) for some \(L\in A\cup B\), that is, iff
        \(w\entails L\) for some \(L\in A\) or \(w\entails L\) for some
        \(L\in B\); by the same lemma this says that \(w\) satisfies \(A\)
        or satisfies \(B\).
  \item[(b)] For finite sets \(\mathcal A\) and \(\mathcal B\) of clauses
        put
        \(\mathcal A\otimes\mathcal B\defeq
        \set{A\cup B:A\in\mathcal A,\ B\in\mathcal B}\), again a finite set
        of clauses.  A valuation \(w\) satisfies \(\mathcal A\otimes\mathcal B\)
        iff it satisfies \(\mathcal A\) or satisfies \(\mathcal B\).  If
        \(w\) satisfies \(\mathcal A\), then for every \(A\in\mathcal A\)
        and \(B\in\mathcal B\) it satisfies \(A\), hence \(A\cup B\) by (a);
        the case where \(w\) satisfies \(\mathcal B\) is the same with the
        roles of \(A\) and \(B\) exchanged.  If \(w\) satisfies neither,
        choose \(A_0\in\mathcal A\) and \(B_0\in\mathcal B\) that \(w\) does
        not satisfy; by (a), \(w\) does not satisfy
        \(A_0\cup B_0\in\mathcal A\otimes\mathcal B\).  In particular
        \(\mathcal A\otimes\varnothing=\varnothing\otimes\mathcal B=\varnothing\),
        in agreement with the statement, since every valuation satisfies
        \(\varnothing\).
\end{enumerate}
The formula set \(\Fm_{\CPL}\) has unique readability
(\cref{thm:prop-unique-readability}), so structural recursion defines, for
every formula \(\varphi\), a pair
\((\mathcal K^{+}(\varphi),\mathcal K^{-}(\varphi))\) of finite sets of
clauses by the following equations, for every \(n\in\N\) and all formulas
\(\psi,\chi\):
\[
\begin{array}{lll}
 \mathcal K^{+}(p_n)\defeq\set{\set{p_n}},
 &\mathcal K^{-}(p_n)\defeq\set{\set{\neg p_n}},\\[2pt]
 \mathcal K^{+}(\bottom)\defeq\set{\varnothing},
 &\mathcal K^{-}(\bottom)\defeq\varnothing,\\[2pt]
 \mathcal K^{+}(\neg\psi)\defeq\mathcal K^{-}(\psi),
 &\mathcal K^{-}(\neg\psi)\defeq\mathcal K^{+}(\psi),\\[2pt]
 \mathcal K^{+}(\psi\land\chi)\defeq
   \mathcal K^{+}(\psi)\cup\mathcal K^{+}(\chi),
 &\mathcal K^{-}(\psi\land\chi)\defeq
   \mathcal K^{-}(\psi)\otimes\mathcal K^{-}(\chi),\\[2pt]
 \mathcal K^{+}(\psi\lor\chi)\defeq
   \mathcal K^{+}(\psi)\otimes\mathcal K^{+}(\chi),
 &\mathcal K^{-}(\psi\lor\chi)\defeq
   \mathcal K^{-}(\psi)\cup\mathcal K^{-}(\chi),\\[2pt]
 \mathcal K^{+}(\psi\to\chi)\defeq
   \mathcal K^{-}(\psi)\otimes\mathcal K^{+}(\chi),
 &\mathcal K^{-}(\psi\to\chi)\defeq
   \mathcal K^{+}(\psi)\cup\mathcal K^{-}(\chi).
\end{array}
\]
Each right-hand side is a pair of finite sets of clauses whenever the
values for \(\psi\) and \(\chi\) are, because unions of finite sets are
finite and \(\mathcal A\otimes\mathcal B\) has at most
\(\abs{\mathcal A}\cdot\abs{\mathcal B}\) members, each the union of two
finite clauses.  The recursion is therefore well defined, and it
terminates on every input: each equation calls the recursion only on the
immediate subformulas of its argument.

\emph{Claim.}  For every formula \(\varphi\):
\begin{enumerate}
  \item[(i)] a valuation \(w\) satisfies \(\mathcal K^{+}(\varphi)\) iff
        \(w\entails\varphi\);
  \item[(ii)] a valuation \(w\) satisfies \(\mathcal K^{-}(\varphi)\) iff
        \(w\not\entails\varphi\);
  \item[(iii)] every literal in a member of
        \(\mathcal K^{+}(\varphi)\cup\mathcal K^{-}(\varphi)\) is \(p_n\) or
        \(\neg p_n\) for a variable \(p_n\) occurring in \(\varphi\).
\end{enumerate}
The proof is by structural induction on \(\varphi\), with (i), (ii) and
(iii) proved together.
\begin{itemize}
  \item \(\varphi=p_n\).  By \cref{lem:prop-clause-satisfaction}(1),
        \(w\) satisfies \(\set{p_n}\) iff \(w\entails p_n\), and satisfies
        \(\set{\neg p_n}\) iff \(w\entails\neg p_n\), that is, iff
        \(w\not\entails p_n\) by the truth clause for \(\neg\).  The only
        literals are \(p_n\) and \(\neg p_n\).
  \item \(\varphi=\bottom\).  No valuation satisfies
        \(\set{\varnothing}\) and no valuation satisfies \(\bottom\); every
        valuation satisfies \(\varnothing\) and every valuation fails to
        satisfy \(\bottom\).  Neither set contains a literal.
  \item \(\varphi=\neg\psi\).  By the induction hypothesis (ii) for
        \(\psi\), \(w\) satisfies \(\mathcal K^{+}(\neg\psi)=\mathcal K^{-}(\psi)\)
        iff \(w\not\entails\psi\), iff \(w\entails\neg\psi\).  By (i) for
        \(\psi\), \(w\) satisfies \(\mathcal K^{-}(\neg\psi)=\mathcal K^{+}(\psi)\)
        iff \(w\entails\psi\), iff \(w\not\entails\neg\psi\).  The variables
        of \(\neg\psi\) are those of \(\psi\).
  \item \(\varphi=(\psi\land\chi)\).  A valuation satisfies a union of two
        sets of clauses iff it satisfies both, so by (i) for \(\psi\) and
        \(\chi\), \(w\) satisfies \(\mathcal K^{+}(\varphi)\) iff
        \(w\entails\psi\) and \(w\entails\chi\), iff \(w\entails\varphi\).
        By (b) and (ii) for \(\psi\) and \(\chi\), \(w\) satisfies
        \(\mathcal K^{-}(\varphi)\) iff \(w\not\entails\psi\) or
        \(w\not\entails\chi\), iff \(w\not\entails\varphi\).
  \item \(\varphi=(\psi\lor\chi)\).  By (b) and (i), \(w\) satisfies
        \(\mathcal K^{+}(\varphi)\) iff \(w\entails\psi\) or
        \(w\entails\chi\), iff \(w\entails\varphi\).  By (ii) and the union
        fact, \(w\) satisfies \(\mathcal K^{-}(\varphi)\) iff
        \(w\not\entails\psi\) and \(w\not\entails\chi\), iff
        \(w\not\entails\varphi\).
  \item \(\varphi=(\psi\to\chi)\).  By (b), (ii) for \(\psi\) and (i) for
        \(\chi\), \(w\) satisfies \(\mathcal K^{+}(\varphi)\) iff
        \(w\not\entails\psi\) or \(w\entails\chi\), which by the truth
        clause for \(\to\) holds iff \(w\entails\varphi\).  By the union
        fact, (i) for \(\psi\) and (ii) for \(\chi\), \(w\) satisfies
        \(\mathcal K^{-}(\varphi)\) iff \(w\entails\psi\) and
        \(w\not\entails\chi\), iff \(w\not\entails\varphi\).
\end{itemize}
In each binary case, every literal of a member of \(\mathcal K^{+}(\varphi)\)
or \(\mathcal K^{-}(\varphi)\) lies in a member of one of
\(\mathcal K^{+}(\psi)\), \(\mathcal K^{-}(\psi)\), \(\mathcal K^{+}(\chi)\)
and \(\mathcal K^{-}(\chi)\), because the members
are taken from these sets directly or are unions \(A\cup B\) of such
members; (iii) for \(\psi\) and \(\chi\) then gives (iii) for \(\varphi\),
whose variables include those of \(\psi\) and \(\chi\).  This proves the
claim.

Now let \(m\defeq\abs{\mathcal K^{+}(\varphi)}\).  If \(m=0\), put
\(C_\varphi\defeq\topv\).  If \(m>0\), list the members of
\(\mathcal K^{+}(\varphi)\) as \(C_1,\ldots,C_m\) in strictly increasing
order of \(\sum_{L\in C}2^{\operatorname{code}(L)}\); this number is
different for different clauses by uniqueness of binary expansions, so
the listing is determined by \(\varphi\).  Put \(C_\varphi\defeq N_1\),
where \(N_m,\ldots,N_1\) are built from \(C_1,\ldots,C_m\) by the recursion
in the definition of conjunctive normal form.  In both cases
\(C_\varphi\) is in conjunctive normal form exactly as defined.  If
\(m>0\), then by \cref{lem:prop-clause-satisfaction}(2) and claim (i), an
arbitrary valuation \(w\) satisfies \(C_\varphi\) iff it satisfies every
\(C_j\), iff it satisfies \(\mathcal K^{+}(\varphi)\), iff
\(w\entails\varphi\).  If \(m=0\), every valuation satisfies
\(\mathcal K^{+}(\varphi)=\varnothing\), so by claim (i) every valuation
satisfies \(\varphi\), and every valuation satisfies
\(\topv=\neg\bottom\) by the truth clauses for \(\neg\) and \(\bottom\).
Hence \(C_\varphi\) is logically equivalent to \(\varphi\) on every
input.

The degenerate inputs are covered by the same construction, with no
separate rule.  Since \(\topv\) is \(\neg\bottom\),
\(\mathcal K^{+}(\topv)=\mathcal K^{-}(\bottom)=\varnothing\) and
\(C_{\topv}=\topv\).  Since \(\mathcal K^{+}(\bottom)=\set{\varnothing}\),
\(m=1\) and
\(C_{\bottom}=\operatorname{form}_{\CPL}(\varnothing)=\bottom\), the
conjunctive normal form whose single clause is empty.  If \(\varphi\) has
no variables, claim (iii) makes every member of \(\mathcal K^{+}(\varphi)\)
empty, so \(\mathcal K^{+}(\varphi)\) is \(\varnothing\) or
\(\set{\varnothing}\); by claim (i) it is \(\varnothing\) when \(\varphi\)
is true under every valuation and \(\set{\varnothing}\) when it is false
under every valuation, so \(C_\varphi\) is \(\topv\) or \(\bottom\)
accordingly.  A valid formula with variables can produce a nonempty set of
clauses each satisfied by every valuation: for \(\varphi=(p_0\lor\neg p_0)\),
\(\mathcal K^{+}(\varphi)=\set{\set{p_0}}\otimes\set{\set{\neg p_0}}
=\set{\set{p_0,\neg p_0}}\) and \(C_\varphi=(p_0\lor\neg p_0)\), a
conjunctive normal form with one clause.  Dually, for the unsatisfiable
\(\varphi=(p_0\land\neg p_0)\), \(\mathcal K^{+}(\varphi)=\set{\set{p_0},\set{\neg p_0}}\)
and \(C_\varphi=(p_0\land\neg p_0)\).
\end{proof}

For example, let \(\varphi=(\neg(p_1\to p_2)\lor p_3)\).  The recursion
gives \(\mathcal K^{+}(p_1\to p_2)=\set{\set{\neg p_1}}\otimes\set{\set{p_2}}
=\set{\set{\neg p_1,p_2}}\) and
\(\mathcal K^{-}(p_1\to p_2)=\set{\set{p_1},\set{\neg p_2}}\), so
\(\mathcal K^{+}(\neg(p_1\to p_2))=\set{\set{p_1},\set{\neg p_2}}\) and
\[
 \mathcal K^{+}(\varphi)
 =\set{\set{p_1},\set{\neg p_2}}\otimes\set{\set{p_3}}
 =\set{\set{p_1,p_3},\set{\neg p_2,p_3}}.
\]
The codes are \(\operatorname{code}(p_1)=2\), \(\operatorname{code}(p_3)=6\)
and \(\operatorname{code}(\neg p_2)=5\), so the clause \(\set{p_1,p_3}\)
(weight \(2^2+2^6=68\)) precedes \(\set{\neg p_2,p_3}\) (weight
\(2^5+2^6=96\)), and
\(C_\varphi=((p_1\lor p_3)\land(\neg p_2\lor p_3))\).  The disjunction and
implication cases take all pairwise unions of clauses, so the number of
clauses can multiply at each such step; the construction preserves logical
equivalence, not size.

For the satisfiability-preserving transformation with fresh variables,
see Section~\ref{sec:p02-tseitin-transformation}: the construction is
Definition~\ref{def:p02-tseitin-construction}, the equisatisfiability and
projection statement is Theorem~\ref{thm:p02-tseitin}, proved in
Annex~\ref{annex:p2.prop.tseitin}, and the size bound is
Theorem~\ref{thm:p02-tseitin-size}, proved in
Annex~\ref{annex:p2.prop.tseitin-size}.

An \(n\)-ary \emph{truth-functional connective} is a connective symbol
interpreted by a fixed function \(\B^n\to\B\).  For a set \(C\) of
truth-functional connectives, the \emph{formulas over \(C\)} are generated
by the propositional variables together with one formation clause per member
of \(C\), and the recursive truth clause of each member applies its
interpreting function to the values of the immediate subformulas.  For a
formula \(\varphi\) whose variables are among \(p_1,\ldots,p_n\), the
\emph{truth function of \(\varphi\) in those variables} is the function
\(f_\varphi:\B^n\to\B\) with
\(f_\varphi(a_1,\ldots,a_n)\defeq\widehat v(\varphi)\) for any valuation
\(v\) with \(v(p_i)=a_i\) for each \(i\); this is well defined by
\cref{lem:prop-coincidence}.

\begin{definition}[Functional completeness]
A set \(C\) of truth-functional connectives is \emph{functionally complete}
iff, for every \(n\in\N\) and every function \(f:\B^n\to\B\), some formula
built from variables \(p_1,\ldots,p_n\) and connectives in \(C\) has truth
function \(f\) in those variables.
\end{definition}

\begin{theorem}[A functionally complete basis]
\label{thm:p02-w02-p02-045-functional-basis}
The set \(\set{\neg,\land,\bottom}\) is functionally complete.
\end{theorem}
\begin{proof}
Given \(f:\B^n\to\B\), form the disjunctive normal form whose conjunctions of
literals correspond exactly to the inputs \(a\in\B^n\) for which
\(f(a)=\True\), as in \cref{thm:prop-normal-forms}.  Replace each occurrence
of \(\alpha\lor\beta\) by
\(\neg(\neg\alpha\land\neg\beta)\).  The empty disjunction is \(\bottom\).
The resulting formula uses only \(\neg,\land,\bottom\) and computes \(f\).
\end{proof}
